% Einheit 2 von 3 – Rückwärts (bank/unabhaengigkeit/e2.jsonl, zusammenbau v0.1)
%% TODO Zeitmarke Einheit 2: Marken-Zeile nicht lesbar
%% TODO Zweigzeile Teil 1: was hier gelernt wird (Satz in Schülersprache aus dem Einheitstitel) – Einheit 2
\einheitenkopf[e2]{Einheit 2 von 3 · Rückwärts}
\zweigzeile{Abitur LK}
\setcounter{aufgabe}{11}

%% TODO Titel: Ich-kann-Satz fehlt in der Bank (Platzhalter: Kettenname) – Nr. 12
%% TODO Anweisung: ein Satz über der ersten Teilaufgabe fehlt in der Bank – Nr. 12
\begin{aufgabe}{Rückwärts}
%% TODO \rechenplatz{n}: Schrittzahl des Grundfalls fehlt in der Bank (nur bei fester Schreibform, 2.3 b)
\begin{teile}
\teil „Die Ereignisse A und B sind stochastisch unabhängig; $P(A) = 0{,}3$ und $P(A \cap B) = 0{,}12$. Berechne $P(B)$.“ Was ist zu tun? \\ \kreuz{prüfen – Produktregel anwenden} \\ \kreuz{nutzen – aus der Unabhängigkeit eine Gleichung gewinnen}
\teil Von den Besuchern eines Museums sind 60\,\% Erwachsene; insgesamt kaufen 35\,\% der Besucher einen Katalog. Welcher Anteil der Erwachsenen müsste einen Katalog kaufen, damit „Erwachsener“ und „Katalogkauf“ stochastisch unabhängig sind? Begründe ohne Rechnung.
\teil Die Vierfeldertafel zeigt Wahrscheinlichkeiten zu den Ereignissen A und B mit einem Parameter p. Für welchen Wert von p sind A und B stochastisch unabhängig? \hfill (Abitur 2021 LK)

\vierfeldertafel{A}{B}{p,3p,4p,0{,}2,0{,}8 - 4p,1 - 4p,p + 0{,}2,0{,}8 - p,1}
\teil Ein Glücksrad hat drei Felder mit den Zahlen 1, 2 und 3; die 1 und die 2 erscheinen je mit der Wahrscheinlichkeit p, die 3 mit $1 - 2p$. Es wird zweimal gedreht. Ereignis E: die erste Drehung zeigt die 3. Ereignis G: beide Drehungen zeigen dieselbe Zahl. E und G sind stochastisch unabhängig. Stelle eine Gleichung auf, aus der p berechnet werden kann, und löse sie \hfill (Abitur 2023 LK).
\teil Die Ereignisse A und B sind stochastisch unabhängig mit $P(A) = 0{,}4$ und $P(A \cap B) = 0{,}1$. Berechne $P(B)$. P(B) = \leerfeld
\teil Ein Gerät ist defekt, wenn mindestens einer der Fehler A oder B auftritt; die beiden Fehler treten stochastisch unabhängig voneinander auf. Ein Gerät ist mit der Wahrscheinlichkeit $0{,}3$ defekt, Fehler A tritt mit der Wahrscheinlichkeit $0{,}2$ auf. Berechne die Wahrscheinlichkeit für Fehler B \hfill (Abitur 2024 LK). P(B) = \leerfeld
\teil Für zwei stochastisch unabhängige Ereignisse A und B gilt $P(B) = P(A) + 0{,}3$ und $P(A \cap \overline{B}) = 0{,}1225$. Bestimme $P(A)$ \hfill (Abitur 2025 LK).
\end{teile}
\end{aufgabe}

%% TODO Titel: Ich-kann-Satz fehlt in der Bank (Platzhalter: Kettenname) – Nr. 13
%% TODO Anweisung: ein Satz über der ersten Teilaufgabe fehlt in der Bank – Nr. 13
\begin{aufgabe}{Rückwärts – Fehler finden, Begründen, Anwendung}
\begin{teile}
\teil A und B sind stochastisch unabhängig mit $P(A) = 0{,}5$ und $P(B) = 0{,}3$. Für $P(A \cap \overline{B})$ rechnet Mia: \rechnung{P(A \cap \overline{B}) &= 0{,}5 - 0{,}3 = 0{,}2} Finde den Fehler.
\teil A und B sind stochastisch unabhängig. Begründe, warum dann auch $P(A \cap \overline{B}) = P(A) \cdot P(\overline{B})$ gilt.
\teil Ein Rauchmelder soll bei Rauch mit mindestens $99\,\%$ Wahrscheinlichkeit anschlagen. Er hat zwei unabhängige Sensoren; der erste erkennt Rauch mit $0{,}9$. Wie zuverlässig muss der zweite Sensor mindestens sein?
\end{teile}
\end{aufgabe}
